\chapter{Orquestração por Ondas e Portões de Qualidade}

\section{Decomposição por Ondas (Wave-based Execution)}
Em vez de permitir que agentes modifiquem múltiplos componentes simultaneamente, o Quantiflux exige a ordenação em Ondas:

\begin{enumerate}
    \item \textbf{Onda 0 (Core Contracts):} Definição de interfaces, DTOs e tipos abstratos.
    \item \textbf{Onda 1 (Adapters \& Persistência):} Implementação de adaptadores de infraestrutura e banco de dados.
    \item \textbf{Onda 2 (Serviços \& Lógica de Negócio):} Construção das regras de negócio consuming os contratos da Onda 0.
    \item \textbf{Onda N (Interface \& Integração Externa):} Camadas de exposição e APIs.
\end{enumerate}

\section{Isolamento de Contexto e Prevenção de Context Bleed}
Cada subagente despachado na Onda $K$ recebe estritamente as especificações da Onda $K$ e os contratos consolidados das Ondas $0 \dots K-1$. Informações de Ondas futuras são omitidas, eliminando alucinações e suposições prematuras.

\section{Portões de Qualidade (Wave Boundary Gates)}
A transição da Onda $K$ para a Onda $K+1$ é bloqueada até que os seguintes critérios sejam atendidos de forma auditável:
\begin{itemize}
    \item Compilação sem erros e aprovação de 100\% da suíte de testes unitários.
    \item Estado do repositório Git verificado e limpo (\texttt{git status}).
    \item Registro de métricas e lições aprendidas consolidado no repositório.
\end{itemize}